\section{Related work}

\label{sec:related-work}

\subsection{Runtime verification and task alignment}

\label{sec:related-work-runtime-verification-and-task-alignment}

VIGIL implements a verify-before-commit protocol for tool-stream injection and introduces the 959-case SIREN benchmark \citep{L01}. It is the closest peer-reviewed structural comparator to M3 because both systems interpose before an effect is committed. DRIFT instead combines dynamic rules with control- and data-level constraints and memory-stream isolation \citep{L02}. The Task Shield verifies whether instructions and tool calls contribute to the user's declared objective \citep{L03}. Together, these systems establish verification, information-flow constraints, and task alignment as complementary runtime strategies.

PlanGuard is a related preprint that isolates planning from untrusted context and applies hierarchical tool and parameter consistency checks \citep{L23}. These approaches provide comparison classes for studying how enforcement structure interacts with the authorization information available to a runtime Guard.

\subsection{Privilege, provenance, and information flow}

\label{sec:related-work-privilege-provenance-and-information-flow}

Progent represents least privilege as symbolic tool-and-argument rules. Its latest arXiv v3 description explicitly distinguishes automatic narrowing from expansion that requires approval \citep{L24}. This overlaps directly with M3's permission-monotonicity and confirmation design. IPIGuard, published at EMNLP 2025, constrains execution through a planned tool-dependency graph \citep{L25}. CaMeL separates trusted control and data flow and uses capabilities to prevent un\-au\-thor\-ized flows \citep{L26}. These systems provide the privilege, dependency, and information-flow foundations for provenance-aware runtime authorization.

FAVA is a July 2026 preprint that translates natural-language tasks into a permission intermediate representation, lowers it to an evidence-backed permission graph, and uses an SMT authorizer before effectful execution \citep{L36}. This formulation makes the relationship between task meaning and executable permission explicit. Our study treats the authorization regime as a paired experimental factor and reports benchmark-specific security, benign utility, confirmation/blocking burden, and operational failure together.

Two 2026 preprints overlap even more closely with provenance-aware authorization. ARGUS constructs an influence-provenance graph for context-dependent decisions \citep{L21}, while AuthGraph compares an authorization graph derived in an isolated context against a parameter-source execution graph \citep{L22}. Both connect action authorization to the sources that influence its parameters. We use these preprints to position the measurement question alongside current provenance-aware designs.

\textit{Auditing Provenance Sensitivity in LLM Agent Action Selection} supplies an especially close 2026 measurement counterpart: it holds task content and policy fixed while changing source authority, then measures whether action selection changes \citep{L35}. Its controlled next-action design differs from M3's multi-step evaluation, but it reinforces the need to separate evidence relevance from authorization. Our experiment extends that question to full tool trajectories, paired runtime Guard conditions, benign utility, \nolinkurl{ASK/DENY} burden, and operational failure.

The contribution is methodological and empirical. We vary authorization information, separate online inputs from evaluator fields, compare policy and full-module conditions, and quantify \nolinkurl{ASK} burden and operational failures. The current paired comparison holds ordinary prompts fixed; historical package comparisons retain their original prompt/runtime assignments.

\subsection{Prompt-level defenses and execution-level outcomes}

\label{sec:related-work-prompt-level-defenses-and-execution-level-outcomes}

Structured queries, spotlighting, instruction hierarchy, and injection detectors act at the prompt, representation, or model-output layer \citep{L07,L30,L31,L32}. These approaches complement execution-level monitoring. The corrected evaluation links proposals, Guard decisions, sandbox receipts, and outcome scoring while reporting completion and failures separately. AgentDojo retains its native security checker, whereas M3-v2 uses observed simulated effects.

\subsection{Benchmarks and adaptive evaluation}

\label{sec:related-work-benchmarks-and-adaptive-evaluation}

AgentDojo supplies dynamic security and utility tasks and is the main public benchmark in this study \citep{L09}. InjecAgent provides a complementary collection of indirect prompt injections for tool-integrated agents \citep{L10}. ASB broadens the scenario and defense taxonomy \citep{L34}. The public expansion retains the benchmarks' unique task identities instead of duplicating benign cases to manufacture a denominator.

\textit{The Attacker Moves Second} motivates defense-aware optimization rather than evaluation against a fixed attack set \citep{L27}. Our historical adaptive panel exposes public \nolinkurl{ALLOW/ASK/DENY} and legacy evaluator-labelled risk feedback, then measures first success on frozen tasks at prespecified query budgets.

PIArena, published at ACL 2026, integrates attacks, defenses, and benchmarks and evaluates feedback-adaptive injection, including cases where the injected task aligns with the intended task \citep{L37}. The AgentDyn preprint instead emphasizes dynamic planning, open-ended tasks, and helpful third-party instructions \citep{L38}. These works motivate testing useful behavior under authorization constraints, not merely counting blocked actions. The ASPI preprint evaluates matched execution and clarification-seeking states and finds increased vulnerability in its clarification setting \citep{L39}. That interaction differs from our controller-only permission diagnostic, which measures permit eligibility without a model clarification turn.

\subsection{Positioning summary}

\label{sec:related-work-positioning-summary}

The active paper studies how authorization scope and runtime configuration shape the measured security--utility frontier. Its paired design connects this question to prior work on verification, least privilege, provenance, and task alignment. Historical studies supply context for prompt/runtime factors and attack-budget sensitivity; the current evaluation supplies corrected authorization and contemporaneous utility measurements.
